\clearpage
\fancyhead[L]{\small BODE. ASTRONOMICAL YEARBOOK.}
\hypertarget{native-630-lower}{}
\begin{center}\small Göttingische gelehrte Anzeigen. No. 5. Pages 41..48. 8 January 1821.\\[3pt]\rule{.72\linewidth}{.3pt}\end{center}

\emph{Astronomical Yearbook for the Year \textup{1823}, with a collection of the latest papers, observations, and reports in the astronomical sciences. Calculated and published by} Dr. J. E. \textsc{Bode}, Royal Astronomer, etc. Berlin 1820. Available from the author and on commission through F. Dümmler. 252 pages in octavo, with one copperplate.

\begingroup\brokenpenalty=10000\baselineskip=13.2pt\relax
Each issue of the Astronomical Yearbook is now doubly welcome to astronomy enthusiasts. Since the astronomical journal edited by Messrs. \textsc{von Lindenau} and \textsc{Bohnenberger} ceased publication, the yearbook has provided the only forum for this science in Germany. The papers included in this volume are as follows. An essay on the physical nature of comets, especially on the comet tail, by Prof. \textsc{Fischer}. It is a pleasure to see here a compilation of what the most reliable observations and more refined views have taught us about this subject. There are, \SourceLeaf{631}however, still several objections to be made to a few individual claims. As is well known, astronomers are still divided over whether comets shine by their own light or by reflected sunlight. In an earlier paper, an astronomer who has made many distinguished contributions to the theory of comets covered virtually everything that can be said on the subject given the present state of our knowledge. The author takes the former view because the light of the smallest fixed stars penetrates even the denser parts of comet tails and, according to some observations, fixed stars are said to have been seen even through the nucleus. Yet he seems to be mistaken in inferring from this that the material of a comet cannot reflect sunlight to any appreciable degree and that the light by which we see comets must necessarily be entirely their own. That observation proves only that the material of a comet transmits far more light than it reflects. But given how very much fainter even the brightest comets appear than planetary disks, there is still no need to invoke light of the comets' own, even if only an extremely small fraction of the sunlight is reflected. In assuming that cometary masses contain substances that gravitate negatively toward the Sun, the author seems to have overlooked that the motions of such comets would then fail to follow \textsc{Kepler}'s laws exactly. Now, while we do not deny that one of \textsc{Kepler}'s laws might indeed be modified somewhat for cometary motion and that this point has never actually been investigated systematically for any comet, one can nevertheless infer from all the calculations made so far that any such modification is likely to be extremely small. --- Observations of the comet of 1819 at the Bogenhausen Observatory near Munich by Tax Councillor \textsc{Soldner}. --- Determinations of geographical positions in East Frisia by Prof. \textsc{Oltmanns}. --- On the longitude of Pisa from astronomical observations, by Prof. \textsc{Wurm}. This valuable investigation shows that the \emph{difference} between the longitude that \textsc{Inghirami} inferred from geodetic measurements and that inferred from astronomical observations cannot really prove anything, given the poor quality of the latter observations. --- By the same astronomer, contributions to geographical determinations of longitude from observations of the solar eclipses of 18 November 1816 and 4 May 1818. --- Astronomical observations at Wilna by Prof. \textsc{Sniadecky}, at Palermo by Mr. \textsc{Cacciatore} (concerning the comets of 1819), and at Prague by Professors \textsc{David} and \textsc{Bittner}. --- A little more about the great comet and its transit across the Sun, by Dr. \textsc{Olbers}. The observation already cited in the astronomical yearbook for 1822 was made by General \textsc{von Lindener}, who on 26 June 1819, in the hour when the comet was certain to be in front of the solar disk, saw the Sun entirely free of spots. It seemed to prove conclusively what was already highly probable, namely that the comet's tenuous mass could intercept far too little sunlight to cause the slightest dimming. Yet this incident shows in a remarkable way how easily errors can arise in such cases when no particularly close attention is paid. Several other observers did in fact see real sunspots on that same day, including Prof. \textsc{Schumacher}, then in Altona, and Prof. \textsc{Brandes} in Breslau. Mr. \textsc{von Lindener}'s observations therefore lose their force as evidence. For if he missed the actual sunspots, he was even less likely to notice the comet, even if it would have been discernible with more powerful instruments or closer attention. Indeed, two other observations might almost lead one to accept that latter possibility, were it not for certain circumstances that make the matter quite uncertain again. On that very day, during the hour when the comet stood in front of the Sun, Dr. \textsc{Gruithuisen} in Munich saw, besides two other ordinary sunspots, a very small one without a distinct boundary that he had not noticed on \SourceLeaf{632} earlier days. Dr. \textsc{Gruithuisen}, however, describes it as a \emph{black} point, and this comet could not, after all, have appeared \emph{like that}. We would therefore be more inclined to regard the phenomenon as an ordinary sunspot that had formed only shortly before and remained invisible in the weaker telescopes of Messrs. \textsc{Schumacher} and \textsc{Brandes}. Prof. \textsc{Wildt} in Hanover also noticed a very diffuse spot on the Sun at the same hour in the morning, but he himself is uncertain of the date, since he reported it only much later, relying solely on his memory. Under these circumstances the main question remains unsettled and may remain so for a long time, since conjunctions of this kind are extremely rare. --- The ephemeris of the Pole Star for all its upper culminations in 1821 will be very welcome to all practical astronomers who do not themselves own the short pamphlet by Messrs. \textsc{Struve} and \textsc{Walbeck} from which it is taken. It might well have deserved an extra page, however, so that the last decimal place would not have had to be omitted. --- Astronomical observations in Vienna by Prof. \textsc{Bürg}, in Prague by Prof. \textsc{Hallaschka}, and in Berlin by the editor. --- On the observed existence of a photosphere of Venus in 1820, by Privy Councillor \textsc{Pastorf}. In April of this year Mr. P. saw Venus (and later Jupiter as well) surrounded by a circular glow with a very sharply defined edge, which he regards as a sphere of light belonging to the planet itself. This phenomenon deserves closer investigation, since an optical illusion can so easily play a part in cases of this kind. In April Mr. \textsc{Pastorf} found that this glow, which resembles the night side of the Moon in earthshine, was 16 arcminutes across. In the phenomenon observed by Mr. \textsc{Pastorf} on 27 April, a small telescopic star entered at the upper part of the \emph{eastern} edge of Venus's sphere of light, remained invisible for \emph{a few minutes}, and then reappeared on the \emph{western side}. The fact that Venus advanced by only 3 arcseconds per minute on that day gives some reason for doubt. We leave it to our readers to judge Mr. \textsc{Pastorf}'s hypothesis that Venus might have a satellite that does not pass beyond the sphere of light and that, as Mr. \textsc{Pastorf} believes, would, if its diameter were only 3 arcseconds, always remain invisible to us in reflected sunlight, except when it appeared from time to time as a dark spot in front of Venus. --- Right ascensions of Saturn and Vesta in 1819 and of Pallas and Mars in 1820, observed with the meridian telescope of the Göttingen Observatory by Court Councillor \textsc{Gauss}. --- Description of the \textsc{Reichenbach} meridian circle installed at the Königsberg Observatory, by Prof. \textsc{Bessel}. --- Observations of the comet of 1819, together with stellar occultations, by Prof. \textsc{Struve}. --- Determination of the obliquity of the ecliptic with a \textsc{Reichenbach} meridian circle at Bogenhausen Observatory near Munich, by Mr. \textsc{Soldner}. The observations used for this purpose also determine the solar semidiameter, giving $16'\,0''93$ at the mean distance. We note on this occasion that 271 wire crossings observed with the \textsc{Reichenbach} meridian telescope at the observatory here from January to July 1820 gave almost exactly the same semidiameter, namely $16'\,1''01$. --- Astronomical observations at Kremsmünster by Mr. \textsc{Derflinger}. --- Observations of Juno, Pallas, Ceres, Mars and Uranus, and of the obliquity of the ecliptic, by Prof. \textsc{Nicolai} in Mannheim. The obliquity, determined with the three-footed \textsc{Reichenbach} repeating circle, comes out $6''$ smaller from the winter solstice than from the summer solstice. This result, together with the stellar declinations determined with the same instrument, confirms once again that each astronomical instrument must be treated as an individual when considering the effect of gravity on its parts. --- On the accuracy of observations with the meridian telescope, by Dr. \textsc{Walbeck}. Here the correct way to treat this \SourceLeaf{633} important subject is taught for the first time and applied to the excellent Dorpat observations. On p.\,186, line\,3 from the bottom, a misprint gives tang $\delta$ instead of sec $\delta$. --- Astronomical observations at Hamburg by Mr. \textsc{Rümker}. --- On the geographical position of Dresden, by Dr. \textsc{Raschig}. --- Astronomical remarks by the Rev. \textsc{Luthmer} in Hanover. --- On the true date of the night battle on the Halys, by Prof. \textsc{Oltmanns}. In this excellent paper, from which we regret to see all citations and notes omitted to save space, it is shown quite convincingly that the much-discussed solar eclipse that occurred during the famous battle can only have been the one on 30 September 609 BC. --- On the orbit of Pons's (\textsc{Enke}'s) comet, together with a calculation of its course on its next return in 1822, by Prof. \textsc{Enke}. The repeated and complete calculation of the perturbations this comet underwent from 1786--1819 gives a most satisfactory fit to the observations from the four apparitions to date. There is, however, a remarkable acceleration of its motion that the calculation cannot yet explain. The three revolutions from 1786--1795 yield an orbital period half a day longer, and the four revolutions from 1805--1819 one half a day shorter, than the three from 1795--1805. If the comet's next appearance yields a similar result, it will be possible to offer plausible conjectures about the cause of this phenomenon. To make it easier to find the comet again in 1822, when it will pass through perihelion on 24 or 25 May, Prof. \textsc{Enke} has very helpfully supplied a convenient ephemeris. Observing it in June and July will present no difficulties in the southern hemisphere. In Europe, however, it will be entirely invisible at that time because of its southerly position, and earlier its extreme faintness will make it very difficult, if not impossible, to detect. To leave nothing untried, we hope that the astronomer who has done so much for the study of this comet will extend the ephemeris, which begins on 25 February, a couple of months further back. Earlier, even though the comet will be fainter still, its high position in the sky on a dark night may offer a little more hope than when the comet is already so near the horizon by the end of twilight. --- For the smaller comet of 1819, Prof. \textsc{Enke} has also found an elliptical orbit with a period of only $5\tfrac{1}{2}$ years, which fits both the Marseille observations and the Milan observations that became known only later extremely well. --- By the same astronomer, the opposition of Vesta in 1819 and an ephemeris for the next apparition of this planet. --- Ephemerides for Juno and Pallas in 1821, by Prof. \textsc{Nicolai} and Mr. \textsc{von Staudt} in Göttingen. --- Astronomical observations in 1820 by Court Councillor \textsc{Gauss}. --- On determining geographical latitude by means of the Pole Star, by Prof. \textsc{Dirksen} in Berlin (analysis of a convenient table provided for this purpose by Prof. \textsc{Schumacher}). --- Further observations of stellar occultations, eclipses of Jupiter's satellites and the great solar eclipse of 7 September 1820 by Prof. \textsc{Rümker} in Hamburg, and of the last of these also by Prof. \textsc{Nicolai} and Mr. \textsc{von Heiligenstein} in Mannheim. --- Of the brief reports at the end of this volume, we also single out those on the new observatory built in Abo and on the arc measurement to be carried out in the Russian Baltic provinces. --- Besides the stellar occultations and the two lunar eclipses of 1823, the copperplate shows the calculated transit of the great comet of 1819 in front of the Sun, the observed geocentric motion of \textsc{Enke}'s comet in 1818 and 1819, and Venus with a spot observed by Mr. \textsc{Pastorff}.
\par\clearpage\endgroup
